\section{Background and Related Work}

\subsection{Quantum Testing as an Oracle Problem}

Metamorphic testing replaces unavailable expected-output oracles with relations between executions of a program~\cite{chen1998metamorphic,barr2015oracle,segura2016survey,chen2018metamorphic}. Prior work reports that relation diversity can alleviate oracle weakness, while still depending on the quality of the chosen relations~\cite{liu2014effectively}. For quantum software, this idea is natural because many behavioral expectations are relational: equivalent circuits should induce the same measurement distribution, a logical qubit permutation should not change results after the output order is restored, and a global phase should not affect measurement probabilities. The difficulty is that these statements hold only under explicit preconditions.

The second difficulty is statistical. A source and follow-up execution can produce different empirical count vectors even when their true distributions are equal, while a real fault may produce a small distribution shift that is invisible under a low shot budget. Treating raw count equality as the oracle creates false alarms, while treating every p-value in isolation ignores the multiple follow-up executions generated per source program. We therefore combine relation validity with a corrected two-sample decision rule for categorical count data.

\subsection{Count Vectors and Output Canonicalization}
\label{sec:count-vectors}

A run of a quantum program produces a \emph{count vector}: a map from measured bitstrings to observed frequencies over $n$ shots. Comparing two runs is comparing two such vectors over a shared support. Two subtleties make this comparison non-trivial.

First, a valid relation often changes the \emph{representation} of the output without changing its logical content. Permuting input qubit roles, remapping a classical register, or reversing a QFT relabels which physical bit carries which logical value. Before any statistical test, QMTester rewrites the follow-up count vector into the source's logical order using a \emph{canonicalization} map $c$ induced by the transformation. A relation is admitted only when $c$ is a documented \emph{bijection} on the measured support: a candidate that drops a register, changes the number of measured bits, introduces a measurement where none existed, or yields a non-bijective map is rejected as a coverage loss, not scored as a test outcome.

Second, Qiskit exposes multi-register outputs as space-separated keys with register-dependent bit order and endianness. Canonicalization first flattens these keys and validates their length, then applies the relation's bijective remap. This is where representation bugs and genuine logical bugs are separated: if two logically equivalent program inputs disagree after a validated bijective remap, the disagreement is a candidate fault rather than a bookkeeping artifact.

\subsection{Non-Unitary Boundaries}
\label{sec:nonunitary}

Mid-circuit measurement, classical conditionals, reset, and other non-unitary operations break the simple ``rewrite-anywhere'' assumption that unitary metamorphic relations rely on. A transformation that is sound for a purely unitary prefix can be invalid once it crosses a measurement or a conditioned reset, because the operation's effect depends on a classical value produced at run time. Program-level relations must therefore declare, per family, whether and how they admit transformations across such boundaries. Several of our relations (e.g.\ ancilla uncompute and teleportation feedback) are specifically about getting these boundaries right, which is why they are expressed over the program's builder semantics rather than over the emitted circuit.

\subsection{Positioning}

We position QMTester by oracle assumption, evidence model, and \emph{testing level} in Table~\ref{tab:related-comparison}. MorphQ is the closest baseline because it also uses metamorphic testing in the Qiskit ecosystem. The decisive distinction is the level at which relations are defined: MorphQ is a \emph{circuit-level} baseline that transforms a generated or given final circuit, whereas QMTester defines relations over \emph{program inputs} and rebuilds the circuit from the follow-up input. This is not a statement that MorphQ is weak. It means MorphQ structurally cannot exploit builder-level semantics: when a defect is already present in the final circuit a builder produced, a circuit-level transformation rewrites a faithful copy of the buggy output and observes no inconsistency. A MorphQ score of $0/7$ on our audited program-level subset is therefore an inherent consequence of the testing level, not an unfair comparison or a tuning artifact.

\begin{table*}[t]
\centering
\caption{Positioning against related quantum testing work. The decisive axis is the \textbf{testing level}: QMTester defines relations over \emph{program/builder inputs}, whereas MorphQ and related metamorphic tools transform a \emph{final circuit}.}
\label{tab:related-comparison}
\scriptsize
\setlength{\tabcolsep}{2.5pt}
\begin{tabular}{@{}p{0.12\textwidth}p{0.10\textwidth}p{0.17\textwidth}p{0.18\textwidth}p{0.30\textwidth}@{}}
\toprule
\textbf{Work} & \textbf{Testing level} & \textbf{Primary focus} & \textbf{Oracle / evidence model} & \textbf{Relationship to QMTester} \\
\midrule
\textbf{QMTester (ours)} & Program / builder & Metamorphic testing of Qiskit program inputs & Program-relation violation via canonicalized count-vector test & --- \\
MorphQ~\cite{paltenghi2023morphq} & Circuit / platform & Metamorphic testing of the Qiskit platform & Generated final-circuit pairs and platform-level failures & Closest baseline; circuit-level, so it cannot reach builder-layer contracts (hence $0$ on the audited program-level subset). \\
Bugs4Q~\cite{miao2023bugs4q} & --- (benchmark) & Reproducible quantum-program bugs & Curated bug/fix pairs & Real-defect substrate; we use an audited program-level subset with documented relations. \\
QDiff~\cite{wang2021qdiff} & Stack & Differential testing across stacks & Cross-stack output discrepancies & Complementary; needs comparable implementations across stacks. \\
QMutPy~\cite{fortunato2022qmutpy} & Circuit / suite & Mutation testing for Qiskit & Existing tests kill or survive mutants & Related only; \emph{not evaluated} here---assumes a test suite, not an oracle-light baseline. \\
QOPS/SB-QOPS~\cite{muqeet2024qops,muqeet2026sbqops} & Circuit & Pauli-string and search-based testing & Measurement-centric Pauli-string oracles & Strong oracle-centered alternative; targets settings with Pauli specifications. \\
QEMI~\cite{luo2026qemi} & Stack & EMI-style stack testing & Dead-code variants should preserve behavior & Complementary metamorphic direction over dead-code variants. \\
Statistical assertions~\cite{huang2019statistical} & Program & Probabilistic program assertions & Developer-specified probability properties & Statistical discipline but requires absolute probability specifications. \\
Projective-measurement testing~\cite{li2020projection,oldfield2025faster} & Circuit & Efficient assertion checking & Projective predicates and reduced measurements & Complementary route to stronger or cheaper relation-specific oracles. \\
\bottomrule
\end{tabular}
\end{table*}


Closest to our setting is the quantum metamorphic-testing line itself. Abreu et al.~\cite{abreu2022metamorphic} and Costa et al.~\cite{costa2022asserting} develop metamorphic relations for oracle and Shor-style quantum programs, while MorphQ~\cite{paltenghi2023morphq} and its reproducibility follow-up MorphQ++~\cite{kitt2024morphqpp} apply metamorphic transformations to the Qiskit quantum platform. These works define relations primarily over the \emph{circuit} (or over an algorithm-level specification) and transform an already-emitted circuit. QMTester differs on two axes that matter for builder-layer defects: (i) its relations are defined over \emph{program/builder inputs} and the follow-up circuit is \emph{rebuilt} from the transformed input rather than rewritten from the source circuit, and (ii) a candidate is admitted only when its output canonicalization is a validated bijection on the measured support. The first axis is what lets a relation see a defect the buggy builder bakes into its circuit; the second is what prevents a representation change from being scored as a fault. We retain MorphQ as the matched circuit-level baseline because it is the most directly runnable tool on the same Qiskit subjects; the contrast with these algorithm/circuit-level relations is about the testing \emph{level}, not tool quality.

The comparison also clarifies what we do not claim. We do not replace proof-oriented formal methods for programs with available specifications, and we do not make Pauli-string or projective assertions unnecessary when such specifications are known. Quantum Hoare logics, relational logics, embedded circuit languages, and circuit-equivalence checkers offer stronger reasoning obligations when their models and specifications apply~\cite{ying2011hoare,unruh2019qrhl,paykin2017qwire,peham2022zx}. We use those ideas as a boundary: QMTester checks empirically whether admitted source/follow-up program pairs behave consistently, but it is not a platform differential-testing campaign, a mutation-adequacy metric, a proof system, or a device-backed study. We also do not evaluate QMutPy in this submission; mutation testing assumes an existing test suite and is not a substitute for the oracle-light, builder-level use case we target.
